\documentclass{beamer}
\usepackage[utf8]{inputenc}
\usepackage{xcolor}

\usetheme{Berkeley}
\usecolortheme{default}

\setbeamertemplate{footline}[frame number]
\beamertemplatenavigationsymbolsempty


\title{Rainbow Paragraphs}
\author[]{IB 514: Scientific Writing}
\date[]{}

\begin{document}

% -------------------------------------------------
\begin{frame}
  \titlepage
\end{frame}


% ============================================================
\begin{frame}{The Rainbow Paragraph Structure}

  \footnotesize
  A \textbf{colour-coded} tool for building and analyzing paragraph structure.
  Each sentence type is assigned a colour to ensure completeness and balance.

  \bigskip
  \begin{description}
    \setlength{\itemsep}{5pt}
    \item[{\color{red}\textbf{Red}} \hfill \textit{Topic Sentence}]
      Introduces the main idea of the paragraph.
    \item[{\color{orange}\textbf{Orange}} \hfill \textit{Context}]
      Provides background information.
    \item[{\color{yellow}\textbf{Yellow}}/{\color{green}\textbf{Green}} \hfill \textit{Evidence / Data / Details}]
      Specific facts, data, or steps in a process.
    \item[{\color{blue}\textbf{Blue}} \hfill \textit{Explanation / Analysis}]
      Explains why or how the evidence supports the topic sentence.
    \item[{\color{purple}\textbf{Purple}} \hfill \textit{Conclusion}]
      Restates the main idea or summarises significance.
  \end{description}

  \bigskip
  \begin{block}{Key rule}
    Every paragraph must begin with a clear
    {\color{red}\textbf{topic sentence}}.
    You do \emph{not} need to use every colour.
  \end{block}

\end{frame}


% ============================================================
\begin{frame}{Why Rainbow Paragraphs?}

  \footnotesize
  \begin{block}{The goal}
    Rainbow paragraphs are a \emph{diagnostic} and \emph{generative} tool ---
    use them to audit existing writing \emph{and} to plan new writing.
  \end{block}

  \medskip
  \textbf{What the colours reveal:}
  \begin{itemize}
    \item \textbf{Missing colours} signal structural gaps
      (e.g.\ no {\color{blue}analysis}, no {\color{purple}conclusion}).
    \item \textbf{Colour imbalance} may indicate an under-supported claim
      or an over-long digression.
    \item \textbf{Ambiguous sentences} --- those that could be two colours ---
      often need to be split or sharpened.
  \end{itemize}

  \medskip
  \textbf{Good defaults to remember:}
  \begin{itemize}
    \item {\color{red}Red} and {\color{purple}purple} each appear \emph{once}.
    \item {\color{orange}Orange}, {\color{olive}yellow/green}, and {\color{blue}blue}
      can each span multiple sentences.
    \item Paragraphs in \textit{Methods} sections naturally use fewer colours.
  \end{itemize}

\end{frame}
% ============================================================
\begin{frame}{Activity}

  \footnotesize
  \textbf{Hopefully:} you brought (1) a written paragraph from your
  \textit{Introduction}, \textit{Results}, or \textit{Discussion},
  and (2) an outline for an unwritten paragraph from either section.

  \medskip
  \textbf{Part 1 --- Colour-code a classmate's paragraph}
  \begin{itemize}
    \item Exchange your written paragraph with a classmate.
    \item Colour-code their paragraph \emph{(up to 15 min)}, then colour-code your own \emph{(up to 10 min)}.
    \item Return paragraphs; compare colour-codings and discuss differences ---
      what made sentences easy or difficult to categorize? \emph{(up to 15 min)}
  \end{itemize}

  \medskip
  \textbf{Part 2 --- Outline an unwritten paragraph}
  \begin{itemize}
    \item Colour-code the outline of your unwritten paragraph and
      elaborate as desired \emph{(up to 10 min)}.
    \item Explain the content and structure of your paragraph to your
      classmate \emph{(10 min each)}.
  \end{itemize}

  \medskip
  We'll use the remaining class period to reflect on the activity as a group.

\end{frame}



\end{document}
